\chapter{Proofs for the constructive spatial RG family}
\label{app:spatial-rg-proofs}

This appendix proves the constructions in
\cref{ch:constructive-spatial-rg}. The geometric inventories, reference
measures, and parameter meanings are those of that chapter. The symbol
\(\mu\) denotes a lattice probability law here, not a Gaussian mean.
\status{DEFINITION}

\section{The source measure and boundary comparison}
\label{sec:rg-spatial-proof-measure}

\paragraph{Proof of \cref{thm:rg-spatial-source-law}.}
For fixed matter signs the coordinate change \(u\mapsto z\) is a
bijection, preserving normalized counting. Its inverse is
\(u_e=\phi_xz_e\phi_{x+e_a}\). Thus the finite raw law is independent
uniform matter times its dressed law in these coordinates. The exterior
\(\eta\) is dressed data. Varying an exterior matter frame while holding
\(\eta\) fixed requires changing its reconstructed raw links as well.

For \(\Delta\subset\Lambda\), conditioning on
\(z_{\Lambda\setminus\Delta}\) cancels the edge factors outside
\(\Delta\) and the plaquettes not meeting \(\Delta\). All remaining factors
are exactly \(W_\Delta\) with the completed exterior. Hence the finite
specification is pointwise consistent. Its one-edge conditional plus
probability is
\begin{equation}
 q_e=\frac{1+\tanh H_e}{2},\qquad
 H_e=h+j\sum_{p\ni e}\prod_{f\in p\setminus\{e\}}z_f,
 \qquad |H_e|\le h+2j.
 \label{eq:rg-spatial-proof-single-row}
\end{equation}
Every conditional sign probability is at least
\(\epsilon_f=(1-\rho)/2>0\). Averaging over any unrevealed variables
preserves this bound, so a specified assignment on \(A\) has probability
at least \(\epsilon_f^{|A|}\).

Let \(n(e,f)\) count plaquettes containing both distinct edges. Each edge
has two plaquettes and three other edges in each, whence
\(\sum_fn(e,f)=6\). Changing one plaquette contribution gives
\begin{equation}
 \frac{|\tanh(s+j)-\tanh(s-j)|}{2}
 =\frac{\sinh(2j)}{\cosh(2s)+\cosh(2j)}\le t.
 \label{eq:rg-spatial-proof-influence}
\end{equation}
Telescoping changes gives influence matrix \(D_{ef}=t n(e,f)\), with
row sums at most \(\alpha=6t<1\).

A random-coordinate heat bath on the finite free set preserves its Gibbs
law by conditional factorization. For \(n=|\Lambda|\), updating each
coordinate in a prescribed order to any target has probability at least
\((\epsilon_f/n)^n\). Therefore every entry of the \(n\)-step transition
matrix has that positive lower bound. Removing its common row mass gives a
strict total-variation contraction, proving uniqueness of the stationary
law. Couple two such chains by the same chosen edge and the same uniform
number in their conditional Bernoulli rows. Cesaro averages of the finite
pair chain have a convergent subsequence; the limit is stationary because
one application changes the average by a difference divided by its length.
Its marginals are the two identified Gibbs laws. If \(d_e\) is its
probability of disagreement, stationarity cancels the unchanged update
probability and gives
\begin{equation}
 d_e\le\sum_{f\in\Lambda}D_{ef}d_f+\beta_e,
 \qquad \beta_e=\sum_{f\notin\Lambda}D_{ef}
                           \mathbf1_{\{\eta_f\ne\zeta_f\}}.
 \label{eq:rg-spatial-proof-stationary-comparison}
\end{equation}
Iteration removes the remainder \(D_\Lambda^n d\), bounded by
\(\alpha^n\). The surviving sum consists of paths making their first exit
from \(\Lambda\) at a differing exterior edge. Dividing each transition
weight by \(\alpha\), and adding a cemetery state for missing row mass,
makes their weights disjoint first-exit probabilities whose sum is at most
one. Every exit path from \(e\) has at least
\(D(e,\Lambda^c)\) steps, so
\(d_e\le\alpha^{D(e,\Lambda^c)}\). A union bound on target disagreements
proves \cref{eq:rg-spatial-boundary-comparison}.

A larger finite-volume restriction is a mixture of the smaller
specifications, by finite consistency. Integrating the same bound over any
two boundary mixtures proves comparison across volumes containing a common
interior. On growing boxes each fixed cylinder is consequently Cauchy,
independently of all exterior choices. Its limit is normalized, consistent,
and has every atom at least \(\epsilon_f^{|A|}\). To construct a measure,
enumerate the countable edges. Split \([0,1)\) into half-open intervals with
the probabilities of the first edge, then split each interval into the
joint probabilities of the next edge. Consistency makes every subdivision
exact. Each point specifies a unique infinite sign sequence; its cylinder
preimages are finite unions of intervals. The pushforward of Lebesgue
measure is a countably additive law with these cylinders. Cylinders form a
pi-system, and the pi-lambda argument makes the extension unique.

For a cylinder function \(f\), the finite update
\(\gamma_\Lambda f\) remains a cylinder function, because the interaction
boundary is finite. Passing finite compatibility to the cylinder limit
gives \(\mu f=\mu\gamma_\Lambda f\). Cylinder determination extends the
identity to bounded measurable functions. For exterior-measurable \(g\),
\(\gamma_\Lambda(gf)=g\gamma_\Lambda f\), so this is the conditional
DLR identity. Any two DLR laws have target marginals that are mixtures of
the same finite specifications. The boundary bound tending to zero makes
all their cylinders equal, hence their measures equal. Translation of the
unique law is another law for the same specification and therefore agrees.
The conditional law in an ordinary box embedded in a full torus has the
same factors and a random dressed exterior. The same comparison proves the
torus cylinder limit while retaining all winding sectors.

For the raw lift, append independent uniform matter signs and apply the
inverse coordinate map. A finite raw window with free matter vertices and
all their incident links converts to independent free frames and the
correct dressed conditional. Such windows are cofinal; conditioning further
proves the raw DLR equations for every finite window. Conversely, a local
gauge flip of one matter vertex and its incident links preserves each raw
conditional density and reference. Any raw Gibbs law is therefore invariant
under every finite product of these flips. In dressed coordinates this
flips arbitrary finite sets of matter signs without changing \(z\).
Averaging those flips proves that every finite matter assignment has
probability \(2^{-|V|}\) independently of every dressed cylinder. Its
dressed marginal has the already identified DLR specification, hence equals
\(\mu_{b,t}\). This proves uniqueness of the raw lift.
\hfill\(\square\)

\section{Constrained parity laws uniformly in blocking depth}
\label{sec:rg-spatial-proof-constrained}

\paragraph{Proof of \cref{thm:rg-spatial-constrained-specification}.}
Partition the fine edges into the retained \(B\)-paths and unused singleton
edges. No plaquette contains two edges of one retained path. A coarse
anchor cell contains two retained paths and \(2B(B-1)\) unused singleton
blocks, hence \(2B^2\) fine edges. A pattern prescribes each retained block
to be unconstrained or to have one specified parity; singleton blocks
remain unconstrained. Use the product of uniform probability references on
these supports, multiply all edge and touching plaquette factors, and
normalize. These are positive finite specifications on their actual
supports. A block row with at most \(B\) edges has every allowed atom at
least
\begin{equation}
 \xi_B=2^{-B}\left(\frac{1-b}{1+b}\right)^B
                   \left(\frac{1-t}{1+t}\right)^{2B}>0.
 \label{eq:rg-spatial-proof-block-positivity}
\end{equation}

Two product sign laws obtained from one positive product law by conditioning
on opposite total parities admit a coupling at Hamming distance exactly
one. For one bit this is immediate. For \(n\) bits, couple the last-bit
marginals by any coupling. If those bits differ, the required prefix
parities agree; couple the prefixes identically. If they agree, apply the
induction coupling to the opposite prefix parities. In either branch the
prefix marginal on each side is its correct product law conditioned on the
required parity, independently of the other branch. This proves the
induction including both marginals.

Changing a single coordinate field from \(+j\) to \(-j\) in a product row
or fixed-parity product row changes that coordinate's marginal by at most
\(t\), by \cref{eq:rg-spatial-proof-influence} with the other-coordinate
masses absorbed into \(s\). Couple it maximally. If it differs in a
fixed-parity row, the remaining prefix parities are opposite, so the
preceding coupling adds exactly one further discrepancy. The expected
Hamming cost is at most \(2t\); a one-bit fixed fiber has zero cost.
Successive tilts can be coupled by composing their conditional coupling
kernels. The endpoint Hamming distance is bounded by the sum of successive
distances, proving the required triangle inequality directly.

For target block \(i\) and exterior fine edge \(e\), let \(n(i,e)\) be
the number of plaquettes meeting both. An exterior flip changes at most
\(n(i,e)\) single-coordinate fields. The preceding coupling gives a row
cost at most \(2t\sum_e n(i,e)\mathbf1_{\{z_e\ne z'_e\}}\).
The column identity, not a block-state row bound, is
\(\sum_i n(i,e)=6\): each of the two plaquettes of \(e\) meets three other
blocks. Incident anchors differ by at most one.

Finite block heat baths have unique stationary laws by the common-mass
argument using \(\xi_B\). A stationary pair coupling exists by the same
Cesaro construction as above. Let \(d_e\) be fine-edge disagreements and
\(D_i=\sum_{e\in i}d_e\). Canceling the update probability separately in
each block gives
\begin{equation}
 D_i\le2t\sum_{e\text{ free}}n(i,e)d_e
      +2t\sum_{e\text{ exterior}}n(i,e)
                         \mathbf1_{\{\eta_e\ne\zeta_e\}}.
 \label{eq:rg-spatial-proof-block-hamming}
\end{equation}
For target anchor set \(A\), multiply by
\(w_i=e^{-\theta\operatorname{dist}_\infty(\operatorname{anchor}(i),A)}\)
and sum. The column identity and \(w_i\le e^\theta w_{\operatorname{block}(e)}\)
when \(n(i,e)>0\) yield
\begin{equation}
 S\le qS+q\sum_{e\in\partial_f\Lambda}w_{\operatorname{block}(e)}
                          \mathbf1_{\{\eta_e\ne\zeta_e\}},\qquad
 S=\sum_{i\text{ free}}w_iD_i.
 \label{eq:rg-spatial-proof-weighted-column}
\end{equation}
Use all blocks anchored in the radius \(R\) neighborhood of \(A\).
Every contributing exterior edge is anchored in its next shell, whose
cardinality is at most \(8|A|(R+1)\), because a square shell around one
cell has exactly \(8(R+1)\) cells. There are \(2B^2\) fine edges per
anchor. Thus the exterior sum is at most
\(16|A|B^2(R+1)e^{-\theta(R+1)}\). Dividing by \(1-q\) and taking a union
bound over target disagreements proves
\cref{eq:rg-spatial-coarse-modulus} for these finite constrained laws.

For every fixed full pattern, the bound makes all finite-cylinder limits
Cauchy on expanding whole-block volumes. Interval subdivision constructs
their measure. Finite compatibility passes to the limit, giving its DLR
law. Any two such laws have cylinders within the same vanishing boundary
bound and hence agree. The argument also compares patterns agreeing in a
large interior; different supports outside it are only exterior forcing.
No conditioning on an infinite event was used.

Finite conditioning on previously free path parities has positive
probability by \cref{eq:rg-spatial-proof-block-positivity}. Restricting its
finite conditional reference is precisely the updated support pattern;
uniqueness identifies the resulting infinite law. In the all-unconstrained
pattern the law is the original fine Gibbs law, because its finite DLR
identities also hold on unions of blocks. Therefore its finite parity
conditionals are these constrained laws. To define a target coarse kernel,
leave its paths unconstrained, prescribe every exterior retained parity,
and take the target marginal of the constructed constrained law.

Conditional on all other block states, the signs along a free path are
independent, with conditional means bounded by \(\rho\). Its parity mean
has magnitude at most \(\rho^B\), also after conditioning on any already
revealed disjoint parities. Revealing the two target parities successively
gives the asserted one-cell lower and upper probabilities. The boundary
comparison makes every kernel continuous and measurable. Conditioning on a
finite set of additionally revealed parities identifies the corresponding
updated pattern; finite Bayes identities therefore give pointwise
specification consistency at every exterior.

For an exterior cylinder event, the defining finite conditional integral
holds under \(\mu_B\). Increase the revealed exterior set; uniform
convergence permits passing the integral to the just constructed kernel.
Cylinder determination extends this conditional identity to the full
exterior sigma algebra. Thus \(\mu_B\) is DLR for this everywhere-defined
specification. Any two coarse DLR laws can be mixed over increasingly
remote exteriors; the same uniform comparison makes their target cylinders
equal. This proves uniqueness. Translation covariance follows by
transporting factors and patterns and invoking constrained uniqueness.
At \(t=0\), independent product-fiber laws give the separate exact control
with coarse edge mean \(b^B\); no formula dividing by \(q\) is needed.
\hfill\(\square\)

\section{The complete canonical interaction}
\label{sec:rg-spatial-proof-potential}

\paragraph{Proof of \cref{thm:rg-spatial-generated-interaction}.}
The one-cell bounds give log-odds amplitude
\(4\operatorname{atanh}(\rho^B)\le A_0\rho^B\). Atom differences are
bounded by total variation, and
\(|\log u-\log v|\le|u-v|/m\) for \(u,v\ge m\), so reference-state
log odds have spatial variation at most
\(D_0B^2(R+1)e^{-\theta(R+1)}\). The canonical supports have diameter at
most \(2r\), and at most \((2r+1)^2\) translates contain the origin.
Their maximal lexicographic cell uniquely specifies the anchor, so no
assignments collide. Consequently
\begin{equation}
 \|\Phi^{(B)}_{A_{x,r}}\|_\infty
 \le\min\{2A_0\rho^B,D_0B^2r e^{-\theta r}\},\qquad r\ge1.
 \label{eq:rg-spatial-proof-shell-minimum}
\end{equation}
The inequality \(\min(u,v)\le\sqrt{uv}\), the bound
\(\sqrt r\le r\), and \(\lambda_v=\theta/8\) reduce the weighted shell
sum to
\(\sqrt{2A_0D_0}B\rho^{B/2}\sum_{r\ge1}r(2r+1)^2z_v^r\).
Differentiating the geometric series three times gives exactly
\(z_v(9+14z_v+z_v^2)/(1-z_v)^4\). This proves the norm bound; its explicit
tail after \(N\) shells replaces that series with
\(\sum_{r>N}r(2r+1)^2z_v^r\). Since
\(u\le e^{\lambda_vu}/\lambda_v\), it also proves the first moment.
Both terms in \(\mathcal E_B\) tend to zero; for an exponential bound,
\(B\rho^{B/4}\le\rho^{1/4}/(1-\rho^{1/4})^2\).

For a configuration differing from \(*\) on finitely many sites, change
those sites in increasing lexicographic order. Positive finite
specification consistency expresses its log probability ratio against the
vacuum as the sum of the corresponding full-past log odds. The absolutely
convergent canonical shells telescope to exactly those odds. Terms anchored
outside the changed set vanish, since their anchor stays at \(*\). Thus
the interaction and specification give identical vacuum-exterior ratios.
Compare two finite-support configurations and cancel their common
exterior to obtain equality for every finite-support exterior. Approximate
an arbitrary exterior by its values in larger boxes and vacuum outside.
Quasilocality gives uniform convergence on the specification side; absolute
summability of all terms touching the finite changed set gives uniform
convergence of interaction energy differences. Positive normalization then
gives equality for every exterior. Infinite total Hamiltonians have not
been used as unnormalized densities.

A continuous specification of a full-support law is unique everywhere:
a difference at one exterior would persist on a cylinder of positive
probability, contradicting almost-sure equality of conditional versions.
The canonical interaction therefore depends on the law. Path composition
and this uniqueness give the stated action on its orbit, even if two
indices happen to describe the same interaction. At the fair product law
all odds are zero, so the interaction is zero and fixed. At \(t=0\) the
singleton interaction is
\begin{equation}
 \Phi_x^{(B)}(s_1,s_2)=-\sum_{a=1}^2
       \log\frac{1+b^Bs_a}{1+b^B},\qquad
 p_B=-2\log(1+b^B),
 \label{eq:rg-spatial-proof-product-vacuum-scalar}
\end{equation}
where \(p_B\) is its normalized-counting pressure. This finite scalar is
not zero despite vacuum normalization. \hfill\(\square\)
\section{The open-ball coefficient map and its law}
\label{sec:rg-spatial-proof-operator}

\paragraph{Proof of \cref{thm:rg-spatial-local-operator}.}
The coefficient space is complete: a norm-Cauchy sequence converges in every
anchored weighted \(\ell^1\) space, its coordinate limits agree on overlaps,
and the uniform Cauchy bound passes to the supremum over anchors.
Translation invariance is a closed coordinate condition. Finite
translation-orbit families of arbitrarily large supports give linearly
independent members, so the positive-radius ball is infinite dimensional.

Condition the fair fine reference on all retained \(a\)-path parities.
Its blocks are independent uniform parity fibers and unused fair
singletons. A character has conditional expectation zero unless its support
is a union of whole paths; a partially occupied path is killed by flipping
one included and one excluded sign, and an occupied unused singleton is
killed by its own flip. A surviving support
\(\operatorname{lift}_a(T)\) has expectation \(\chi_T\).
For an ordered list of source occurrences \(S_1,\ldots,S_n\), repetitions
allowed, every product of moments in the joint cumulant is therefore either
zero or the same character corresponding to
\(S_1\mathbin\triangle\cdots\mathbin\triangle S_n\). Denote its constant
coefficient by \(c_a(S_1,\ldots,S_n)\). The degree-\(n\) map is
\begin{equation}
 P_{a,n}(J_1,\ldots,J_n)_T
 =\frac1{n!}\sum_{S_1\triangle\cdots\triangle S_n=\operatorname{lift}_a(T)}
 c_a(S_1,\ldots,S_n)\prod_{i=1}^n(J_i)_{S_i},\qquad T\ne\varnothing.
 \label{eq:rg-spatial-proof-cumulant-map}
\end{equation}
For finite source inventories this is differentiation of the finite
conditional log partition, including its \(1/n!\) and repeated occurrences.

Join two occurrence vertices when their supports meet a common reference
block. Disconnected groups then depend on independent reference variables.
The bounded dependency-graph cumulant inequality is
\(|c_a|\le2^{n-1}\operatorname{ST}(G)\), where \(\operatorname{ST}\)
counts spanning trees of this occurrence graph
\citep{feray2016modphi}. Its hypotheses hold uniformly in every parity
assignment: the variables are sign characters of magnitude one, and the
independence is at the block level. This is the finite inequality used
here; convergence and the infinite-law identification are proved below.

Write \(s_i=|S_i|\) and \(d_i=\operatorname{diam}S_i\). For nonempty
output \(T\), disjoint paths give \(a|T|\le\sum_i s_i\). Along an
occurrence tree, travel within a source at cost at most \(d_i\), and
between two source edges in one reference block at cost at most \(a-1\).
The first fine edges of two output paths realizing its diameter then give
\begin{equation}
 a\operatorname{diam}T\le\sum_i d_i+(a-1)(n-1).
 \label{eq:rg-spatial-proof-tree-diameter}
\end{equation}
Fix an output edge and its first fine edge. At least one occurrence contains
that fine edge, so sum over \(n\) possible roots. Pay \(|T|\) by
\(a^{-1}\sum_i s_i\), summing over \(n\) marked occurrences. A child
support touching a parent reference block contains an edge in the union of
blocks met by the parent, at most \(a s_i\) edges. Expanding this witness
condition and summing from leaves costs \(a\) per tree edge and one parent
support-size factor per child.

For an occurrence with \(k_i\) children, the weighted anchored sum is
bounded by its input norm times \(k_i!\delta_a^{-k_i}\). Indeed the
unmarked residual factor is
\(s_i^{k_i-1}e^{-\delta_as_i}\), and the marked one is
\(s_i^{k_i}e^{-\delta_as_i}\); both are at most
\(k_i!\delta_a^{-k_i}\), using
\((\delta_as_i)^{k_i}/k_i!\le e^{\delta_as_i}\) and \(s_i\ge1\).
The unused spatial weight only decreases the bound. Summing
\(\prod_i k_i!\) over labeled trees with fixed root gives
\((n-1)!\operatorname{Cat}_{n-1}\): ordering the children counts plane
rooted trees, and the other \(n-1\) labels are free. Balanced-parenthesis
encoding gives \(\operatorname{Cat}_{n-1}\le4^{n-1}\). Combining the
cumulant tree inequality, roots, mark, \(1/n!\), and
\cref{eq:rg-spatial-proof-tree-diameter} proves
\begin{equation}
 \|P_{a,n}(J_1,\ldots,J_n)\|
 \le\frac n a K_a^{n-1}\prod_{i=1}^n\|J_i\|.
 \label{eq:rg-spatial-proof-multilinear-majorant}
\end{equation}
All rearrangements were of positive majorants, so they are justified also
for infinite source inventories.

Define \(\mathcal R_aJ=\sum_{n\ge1}P_{a,n}(J,\ldots,J)\) for
\(K_a\|J\|<1\). The majorant gives locally uniform convergence in the
Banach norm and locally uniform convergence of
\(D\mathcal R_a(J)H=\sum_{n\ge1}nP_{a,n}(H,J,\ldots,J)\) in operator
norm. The multilinear increment expansion leaves a uniformly summable
remainder with at least two increments; hence this is the actual
Fr\'echet derivative, and its locally uniform series is continuous.
The first-order survival rule gives \(L_a\). The lift is injective, has
size \(a|T|\), contains the first edge of every output path, and has
diameter at least \(a\operatorname{diam}T\), so comparing anchored weights
gives \(\|L_a\|\le\ell_a\). For \(x\le1/2\),
\(\sum_{n\ge2}nx^{n-2}\le6\) and
\(\sum_{n\ge2}n^2x^{n-2}\le22\), by differentiating the geometric
series. These prove the quadratic bound with \(6K_a/a\) and the derivative
bound with \(22K_a/a=2M_a\), hence the stated bounds. The chosen radius
makes the radial contraction at most \((1+\ell_a)/2\).

It remains to identify an actual normalized law. For any real input in the
ball, its one-edge half-log-odds are
\(\sum_{S\ni e}J_S\chi_{S\setminus\{e\}}\), absolutely and uniformly
convergent. Flipping \(f\) changes the Bernoulli row by at most
\(\sum_{S\ni e,f}|J_S|\). Its influence row sums are at most
\(\|J\|<1\). Finite heat-bath coupling gives the resolvent comparison
\(d\le\sum_{m\ge0}D^m\beta\). For a fixed target, the contribution
from a remote boundary tends to zero: each fixed power has summable rows,
and the remaining powers have a common geometric tail. Finite-cylinder
subsequences exist by finite-state compactness; uniform convergence of
finite local conditional energies passes their DLR equations to the
limit. The comparison then makes the limit unique and makes the entire
exhaustion converge.

Take finite edge boxes \(E_m\) and source truncations
\(J^{[m]}_S=J_S\mathbf1_{\{S\subset E_m\}}\). Each is a finite inventory
and has no larger norm. No global norm convergence of these inhomogeneous
truncations is asserted. Every fixed source tuple is eventually present;
the rooted majorant therefore gives, for every fixed output edge, weighted
incident convergence of \(\mathcal R_aJ^{[m]}\) to
\(\mathcal R_aJ\). In the finite inventory the convergent scalar and
nonempty Taylor series exponentiate to the exact conditional moment
generating function near zero. The identity of the convergent analytic
series extends over its real convergence interval, and that function is
strictly positive there. Its real logarithm therefore has precisely these
coefficients throughout the admitted ball. The scalar cancels only when
the finite density is normalized.

The truncated fine law is the free-boundary Gibbs law on \(E_m\), with
fair signs outside. A path crossing its boundary retains those outside
fair signs; incomplete characters are canceled rather than discarded by
an altered channel. Its path push has exactly the finite coefficient
expansion above. Incident convergence makes every finite coarse energy
difference, hence every normalized conditional row, converge uniformly in
the exterior. The fine laws converge locally to their unique law and the
path map preserves convergence of cylinder observables. Passing the finite
coarse DLR integral to this limit proves that the path-pushed law is Gibbs
for \(\mathcal R_aJ\). Its norm lies in the ball, so it is the unique
coarse Gibbs law by the same comparison.

For completeness, every prescribed parity exterior has a continuous
conditional version. In the whole-reference-block representation,
changing one outside block gives an influence row sum bounded by
\(2a\|J\|<1\), since a block has at most \(a\) fine edges. This estimate
holds on every fixed-parity or unconstrained support, with all crossing
sources present. Its summable resolvent constructs unique laws for every
pattern and gives local continuity when remote patterns change. Free
retained parity probabilities are at least
\(e^{-2a\|J\|}/2>0\). Finite parity conditioning corresponds exactly to
these product-fiber supports. Testing the uniform limit of its finite
conditional rows against exterior cylinders identifies it as a conditional
version of the path-pushed law. The coefficient Gibbs kernel is also
continuous. Since the law has full cylinder support, the two versions
agree everywhere.

Finally, integrating the conditional half-log-odds at \(e\) against the
fair exterior character \(\chi_{S\setminus\{e\}}\) recovers \(J_S\)
by orthogonality and uniform convergence. This proves uniqueness of every
nonempty coefficient. Path composition, unique laws, and this extraction
give \(\mathcal R_2\mathcal R_3=\mathcal R_3\mathcal R_2=\mathcal R_6\).
The same reasoning applies to any finite word with its directly defined
path channel. Zero is the fair product interaction and is fixed. The
original field and plaquette source has norm
\(e^\kappa h+8e^{4\kappa+\lambda}j\), so the ball includes a nonempty
region with both positive couplings. \hfill\(\square\)

\section{Complete finite Taylor calculations and their errors}
\label{sec:rg-spatial-proof-finite-amplitude}

\paragraph{Proof of \cref{thm:rg-spatial-finite-amplitude}.}
Represent each nonempty support once modulo integer translation, by moving
its lexicographically first edge basepoint to zero and retaining its axis.
If its shape is \(S\), size \(s\), diameter \(d\), and it contains
\(n_i(S)\) axis-\(i\) edges, exactly \(n_i(S)\) distinct translates
contain a specified axis-\(i\) edge. A nonempty finite set has no nonzero
translation period. Therefore a finite library's exact norm is
\begin{equation}
 \|J\|=\max_{i=1,2}\sum_S n_i(S)|S|
                  e^{\kappa|S|+\lambda\operatorname{diam}S}|J_S|.
 \label{eq:rg-spatial-proof-finite-orbit-norm}
\end{equation}
Duplicates are consolidated by exact addition before this formula is used.
Only translations are identified; rotations and reflections are not.

At order \(n\), let \(d_{\max}\) and \(s_{\max}\) be the largest input
shape diameter and size. A connected source-occurrence list contributing
an output with first edge at the origin stays within fine basepoint radius
\(R_n=nd_{\max}+(n-1)(a-1)\), by the occurrence-tree path argument.
Thus every contributing literal instance belongs to the finite catalog
of all input shapes translated by \([-R_n,R_n]^2\cap\mathbb Z^2\).
Reference-block ownership uses Euclidean division: a horizontal edge is
retained exactly when its transverse coordinate is divisible by \(a\),
and its owner uses the floor of its parallel coordinate divided by \(a\);
vertical edges have the rotated rule. Negative coordinates follow the same
floor rule. Unused singleton owners have a distinct type. A support is a
union of complete paths exactly when every owner is retained and has all
\(a\) distinct members.

For a set \(I\) of occurrence labels, let \(m(I)\) be one if their XOR
support is a union of complete paths, including the empty support, and zero
otherwise. At all positive parities this is precisely its conditional
moment. With \(r\) the first occurrence label in a nonempty \(I\), the
cumulant obeys
\begin{equation}
 c(I)=m(I)-\sum_{\substack{D\subsetneq I\\r\in D}}
                           c(D)m(I\setminus D).
 \label{eq:rg-spatial-proof-cumulant-recurrence}
\end{equation}
The partition block containing \(r\) uniquely decomposes the
moment--cumulant formula; summing all other blocks gives the second moment
factor. This proves the recurrence and its singleton base. Occurrence
labels remain distinct even when their supports repeat. Computing all
subsets takes finitely many exact integer operations; fewer than
\(3^n\) subset-pair contributions suffice.

For each order and each output first axis, enumerate every size-\(n\)
multiset of literal catalog instances once. Keep connected occurrence
graphs with nonempty complete-path XOR support whose output is already
canonically anchored on that axis. A multiset with multiplicities
\(m_\tau\) contributes exactly
\(c(I)\prod_\tau J_\tau^{m_\tau}/\prod_\tau m_\tau!\): its
\(n!/\prod m_\tau!\) ordered occurrences cancel the \(n!\) in
\cref{eq:rg-spatial-proof-cumulant-map}. Root choices and spanning trees
are bounds, not extra multiplicities. Disconnected occurrences have zero
cumulant by block independence. The radius argument loses no connected
term. Output canonicalization is a selection of the already canonical
output, not accumulation of all translated copies. Consolidation over all
accepted terms gives the complete coefficient at every support through the
specified order. Surviving supports satisfy
\(a|T|\le ns_{\max}\) and
\(a\operatorname{diam}T\le nd_{\max}+(a-1)(n-1)\), so the output is
finite. Every finite catalog, multiset list, recurrence, and rational
operation terminates at fixed order.

Summing \cref{eq:rg-spatial-proof-multilinear-majorant} over orders above
\(N\) gives \cref{eq:rg-spatial-taylor-tail} by differentiating the
geometric tail \(x^{N+1}/(1-x)\). Its positive-series representation
makes it monotone in the norm and \(K_a\), so certified upper bounds for
both are valid despite the minus sign in the closed numerator. It tends to
zero for \(x<1\). The zero input gives the zero table and zero tail.

All transcendental constants admit convergent rational enclosures. For a
nonnegative rational \(x\), put \(t_0=1\),
\(t_{k+1}=t_kx/(k+1)\), \(s_m=\sum_{k=0}^mt_k\). Once
\(v=x/(m+2)<1\), positivity and decreasing successive term ratios give
\begin{equation}
 s_m\le e^x\le s_m+\frac{t_{m+1}}{1-v}.
 \label{eq:rg-spatial-proof-exp-enclosure}
\end{equation}
The width tends to zero. Reciprocal endpoints enclose \(e^{-x}\).
Applying these termwise to \cref{eq:rg-spatial-proof-finite-orbit-norm}
and \cref{eq:rg-spatial-local-radius} gives convergent upper norm and lower
radius bounds. Every strict admission or positive tolerance gap is
therefore decided at some finite precision. The strict image margin and
vanishing Taylor tail eventually admit the approximate image as well.

For one certified \(r<\varepsilon\), the derivative estimate along a
segment gives Lipschitz constant \(L_a(r)<1\), while the value estimate
gives radial factor \(q_a(r)=\ell_a+M_ar<1\). The triangle inequality
separates one step's Taylor error from its input error and gives the
stated recurrence. Choosing every tail below both a prescribed small
fraction of the total tolerance and \((1-q_a(r))r/2\) keeps all
approximate states in the ball and achieves the requested error for a
finite word. Two differently truncated direct and staged tables need not
be equal; their difference is bounded by the sum of their errors relative
to the same exact composed operator. The empty character was excluded only
from the coefficient table. For example, an unused fine field has a
nontrivial \(\cosh\) normalizer although its retained interaction is zero;
no evidence value can be inferred from that table. \hfill\(\square\)
\section{A constructive basin-entry theorem in the prescribed norm}
\label{sec:rg-spatial-proof-basin}

\paragraph{Proof of \cref{thm:rg-spatial-basin-entrance}.}
The proof uses the specified compact box and retains the prescribed
\(\kappa,\lambda\). For \(B\ge3\), conditionally integrate all but the
first and last edges of each retained path. No plaquette meets two
eliminated edges, and the other three edges in its plaquette are outside
every retained path. These interiors are therefore conditionally
independent. For a horizontal path at the origin its conditional parity
mean is
\begin{equation}
 g_p=z_{(0,0,1)}z_{(B-1,0,1)}
        \prod_{s=1}^{B-2}f(u_s,v_s),\qquad
 f(u,v)=\tanh(h+j(u+v)),
 \label{eq:rg-spatial-proof-basin-strip}
\end{equation}
where \(u_s,v_s\) are the three-edge characters in the adjacent
plaquettes. Its strip consists of its endpoints, the horizontal edges at
\((s,\pm1)\), \(1\le s\le B-2\), and the vertical edges at
\((s,0),(s,-1)\), \(1\le s\le B-1\). It has \(4B-4\) edges and
diameter at most \(B\). These strips and their rotations are pairwise
disjoint. For horizontal edges, horizontal-strip interiors have first
coordinate residues \(1,\ldots,B-2\), whereas vertical-strip interiors
have residues \(0,B-1\); the transverse residues identify the unique
adjacent grid line for \(B\ge3\). Retained endpoints belong only to their
own path. Rotation proves the other orientation. This also proves
conditional independence of all eliminated interiors.

Set \(v=2t/(1+t^2)\). The four Walsh coefficients of \(f\) are
\(a_0,a_1,a_1,a_{12}\), with
\begin{equation}
 a_0-a_{12}=b,\qquad
 2a_1=\frac{v(1-b^2)}{1-b^2v^2},\qquad
 a_{12}=-\frac{bv^2(1-b^2)}{2(1-b^2v^2)}.
 \label{eq:rg-spatial-proof-strip-coefficients}
\end{equation}
Here \(a_0,a_1\ge0\), \(a_{12}\le0\), so the unweighted coefficient
sum is \(R(b,v)=b+v(1-b^2)/(1-b^2v^2)\). It increases in \(v\), and
\(\partial_bR\ge1-2v/(1-v^2)>0\) for
\(v\le24/145\), the value at \(t=1/12\). Thus its upper-endpoint value
\(R_*\) bounds the compact box, with the explicit slack
\(1-R_*\ge(13489/20449)(1-b_+)>0\). Put
\begin{equation}
 \Delta=1-R_*,\qquad \alpha=\Delta/24,\qquad
 \sigma=\frac{R_*}{1-\Delta/4}<1,\qquad
 G_B=\frac{\sigma^{B-2}}{1-2\alpha}.
 \label{eq:rg-spatial-proof-basin-activity}
\end{equation}
The fine Walsh norm with weight \(e^{\alpha|S|}\), including its scalar,
is submultiplicative under symmetric difference. Since \(u,v\) have
three edges each,
\(\|f\|_\alpha\le e^{6\alpha}R_*\le\sigma\) and
\(\|g_p\|_\alpha\le G_B\), using
\(e^x\le1/(1-x)\) for \(0\le x<1\), including the two endpoints.

For a finite path set \(P\), the actual parity density relative to fair
counting is
\begin{equation}
 f_P(y)=\E_{\mu_{b,t}}\prod_{p\in P}(1+y_pg_p).
 \label{eq:rg-spatial-proof-soft-parity-density}
\end{equation}
Conditional independence of the eliminated interiors proves this identity;
averaging over \(y\) gives one. The auxiliary fine law tilted by this
product has the correct remaining-sign marginal for these observations.
It is not the hard-parity posterior on all original fine variables. This
typing matters for its normalized conditional rows below.

For a root path \(p_0\) with fine anchor \(x_0\), let
\(D_{p_0}(S,T)\) be the maximum fine-coordinate distance from \(x_0\)
of any fine-edge anchor in \(S\), or of \(B\) times any coarse-edge
anchor in \(T\). An empty maximum is zero. Set \(\beta=\kappa+1\) and
\(\zeta=2\lambda/B\), and complete the joint fine/coarse Walsh algebra
in
\begin{equation}
 \|F\|_{p_0}=\sum_{S,T}|F_{S,T}|
     e^{\alpha|S|+\beta|T|+\zeta D_{p_0}(S,T)}.
 \label{eq:rg-spatial-proof-rooted-algebra}
\end{equation}
Fine signs and coarse labels are distinct variables. Cardinalities and
rooted radii are subadditive under multiplication, proving the Banach
algebra inequality. For a local update at a fine edge use the same norm
centered there. Removing that edge lowers the fine weight by
\(e^{-\alpha}\); replacing it by a local polynomial increases the rooted
radius by at most that polynomial's local radius.

Put \(p=(1+\sigma)/2\), and increase \(B\) until
\(2\lambda/B<1\) and
\(\sigma/(1-2\lambda/B)\le p<1\). The original conditional mean at an
edge has local norm at most \(p\), because it reaches distance at most one.
Choose rational upper bounds \(A_\beta\ge e^\beta\),
\(A_{2\lambda}\ge e^{2\lambda}\), and
\(A_\kappa\ge e^\kappa\), and set
\(H_B=A_\beta A_{2\lambda}G_B\). Each observed strip has radius at most
\(B\) from any of its members, hence local activity at most \(H_B\).

A fine edge belongs to at most one observed strip. If it belongs to that
strip, write \(g_p=g_0+z_eg_1\) and let \(f_e\) be its original
conditional mean. Direct normalization of the two probabilities after
multiplying by the complete observed factor gives
\begin{equation}
 F_e=\frac{f_e+y_p(f_eg_0+g_1)}{1+y_p(g_0+f_eg_1)}
     =f_e+\frac{y_p(1-f_e^2)g_1}{1+y_p(g_0+f_eg_1)}.
 \label{eq:rg-spatial-proof-normalized-tilted-row}
\end{equation}
If no observed strip touches the edge, use \(F_e=f_e\); omit the target
strip when computing its own conditional mean. Both \(y_pg_0\) and
\(y_pg_1\) have local norm at most \(H_B\). Increase \(B\) until
\begin{equation}
 H_B\le\min\left\{\frac1{2(1+p)},
                       \frac{1-p}{4(1+p^2)}\right\}.
 \label{eq:rg-spatial-proof-basin-row-gate}
\end{equation}
The denominator perturbation then has norm at most \(1/2\), so its
geometric inverse has norm at most two. The correction has norm at most
\(2(1+p^2)H_B\le(1-p)/2\). Thus every normalized conditional mean has
local norm at most \(r=(1+p)/2<1\). The denominator and its even
contribution have both been retained.

Let \(\Pi\) project onto terms with empty fine support. For nonempty
\(S\), define
\begin{equation}
 \mathcal A(\chi_Sy_T)=\frac1{|S|}\sum_{e\in S}
              \chi_{S\setminus\{e\}}y_TF_e,\qquad
 N=(1-\Pi)\mathcal A,
 \label{eq:rg-spatial-proof-absorbing-operator}
\end{equation}
and let \(\mathcal A\) vanish on fine-constant terms. Removing the fine
edge, using its local row norm, and averaging give
\(\|\mathcal A(1-\Pi)\|\le e^{-\alpha}r\le r\). The average is
necessary to avoid a support-size factor. Conditional expectation under
the actual normalized finite tilted law satisfies
\(\E\chi_Sy_T=\E\mathcal A(\chi_Sy_T)\). Absorbing the fine-constant
terms once and iterating proves
\begin{equation}
 \E F=\Pi F+\sum_{k\ge1}\Pi\mathcal A N^{k-1}(1-\Pi)F.
 \label{eq:rg-spatial-proof-absorbing-series}
\end{equation}
The unabsorbed remainder is bounded geometrically in the same norm. At any
fixed substitution order only finitely many fine variables and coarse
labels are reached; the coefficients stabilize once the finite inventories
contain that region. Fixing exterior signs decreases norms. Finite weights
include every crossing original and observed factor. The uniform
geometric remainder therefore proves boundary-independent continuous
limits for arbitrary full observed patterns.

For a finite target set \(A\), observing its complement gives the
conditional density
\(\E_{\mu^{y,A^c}}\prod_{e\in A}(1+x_eg_e)\), relative to fair
counting on \(A\). It is normalized, and lies between
\((1-G_B)^{|A|}\) and \((1+G_B)^{|A|}\). Increasing finite observed
subsets gives this formula uniformly for every exterior by the preceding
stabilization and tail. Finite Bayes consistency passes to these limits.
Testing their conditional integrals against exterior cylinders identifies
them as the actual pushed law's DLR kernels. This supplies the full finite
kernels, not only one mean or an almost-sure version.

For a single target, start \cref{eq:rg-spatial-proof-absorbing-series} at
\(g_{p_0}\), whose rooted norm is at most \(A_{2\lambda}G_B\). Its
conditional mean \(m_{p_0}\) therefore has norm at most
\(M_B=A_{2\lambda}G_B/(1-r)\). Increase \(B\) until \(M_B\le1/2\).
The norm-convergent series
\(u_{p_0}=\operatorname{atanh}m_{p_0}\) gives its complete half-log-odds
and bound \(\|u_{p_0}\|\le U_B=M_B/(1-M_B^2)\). Positive two-edge
conditional kernels satisfy the conditional square identity
\begin{equation}
 u_e(y^{f,+})-u_e(y^{f,-})
 =u_f(y^{e,+})-u_f(y^{e,-}).
 \label{eq:rg-spatial-proof-conditional-square}
\end{equation}
Both sides are half the same log cross ratio. Taking Walsh coefficients,
justified by absolute convergence, shows that the coefficient of
\(\chi_{T\setminus\{e\}}\) in \(u_e\) is independent of the marked
edge \(e\in T\). Define it to be \(J_T^{(B)}\). It reproduces every
single-edge ratio; connectedness of the finite sign cube and normalization
then reproduce all the conditional kernels. Translation covariance gives
one canonical orbit representative.

For a size-\(n\) orbit its translated incidence and norm prefactor are at
most \(n^2\). The inequality
\(n^2e^{\kappa n}\le2e^\kappa e^{(\kappa+1)(n-1)}\) follows from
\(\sup_{x\ge1}x^2e^{-(x-1)}=4/e<2\). Because the root belongs to
\(T\), its diameter is at most twice its rooted radius. Selecting only
supports whose literal first edge is the root and summing the two root
orientations therefore gives
\begin{equation}
 \|J^{(B)}\|_{\kappa,\lambda}
 \le\mathcal B_B:=4A_\kappa\frac{M_B}{1-M_B^2}.
 \label{eq:rg-spatial-proof-basin-norm}
\end{equation}
This estimate counts all supports and translated incidences in the
original prescribed weights. The fine weight \(\alpha\) was auxiliary,
and \(\lambda\) has not been replaced by a smaller output weight.

Every displayed gate is eventually true and remains true for larger
\(B\): the base gate improves with \(B\), and \(G_B\) decreases
exponentially. Rational exponential enclosures as in
\cref{eq:rg-spatial-proof-exp-enclosure} give a positive lower bound
\(\varepsilon_0<\varepsilon\). A finite rational search for \(B_0\)
ensuring the gates and \(\mathcal B_B<\varepsilon_0/16\) terminates,
uniformly over the compact box. The local RG theorem now supplies
uniqueness in its ball. Extraction of conditional coefficients and exact
path composition identify \(\mathcal R_aJ^{(B)}=J^{(aB)}\).

To obtain a finite rational interaction, keep absorbed substitution orders
through \(K\), obtaining \(m_K\), and keep the odd powers of
\(\operatorname{atanh}m_K\) through \(2L-1\). Every local row in
\cref{eq:rg-spatial-proof-normalized-tilted-row} is a rational function of
finitely many signs when \(b,t\) are rational. Its denominator is
positive. Evaluate it on the complete cube and take its finite Walsh
transform, retaining every coefficient. One may reduce the cube exactly
to a basis for the \(\mathbb F_2\)-span of its character supports: the
basis-character map is surjective by row rank; the function is constant
on its fibers; every outside-span coefficient vanishes by pairing with a
kernel element that flips that character. Thus the complete transform on
the basis cube, lifted to literal supports, is the entire raw-coordinate
function. The weights still use literal cardinalities and distances, not
the number of basis generators. The local strip row's character rank is
at most \(2B\), from its \(2(B-2)\) three-edge characters, endpoint
character, possibly removed edge, and two original row characters.

All finite substitution orders and powers have finite inventories. The
conservative mean remainder is \(d_K=M_Br^K\). Both exact and truncated
absorbed sums have norm at most \(M_B\). Telescoping the odd powers gives
an \(\operatorname{atanh}\) Lipschitz bound \((1-M_B^2)^{-1}\), and its
omitted powers give the complete error
\begin{equation}
 \|\widehat J-J^{(B)}\|_{\kappa,\lambda}
 \le4A_\kappa\left[
       \frac{d_K}{1-M_B^2}
       +\frac{M_B^{2L+1}}{(2L+1)(1-M_B^2)}\right].
 \label{eq:rg-spatial-proof-basin-oracle-error}
\end{equation}
The incidence argument applies also to omitted coefficients and to
approximate rows that need not obey exact square identities. Canonical
selection gives one valid approximate interaction. After \(K\)
substitutions its labels lie within coarse radius \(K+1\); the full
\(\operatorname{atanh}\) of that finite-variable mean introduces no new
labels. Hence supports of diameter greater than \(2(K+1)\) have total
norm at most \(4A_\kappa d_K/(1-M_B^2)\). The finite odd polynomial pays
its separate power tail in
\cref{eq:rg-spatial-proof-basin-oracle-error}.

Increasing \(K,L\) makes the error arbitrarily small, and each complete
rational transform, substitution, and product terminates with sufficient
finite resources. A bounded incomplete table does not have this
conclusion. For a finite RG word, choose the input error small enough that
both the true interaction and its finite approximation lie in one
certified radius. Comparing their exact trajectories by the Lipschitz
bounds gives \(\delta\prod_kL_{a_k}(r)\); adding the separate Taylor
word error proves the stated total. The argument retains normalized fine
tilts throughout but computes no discarded original scalar partition.
\hfill\(\square\)
\section{Higher response and the actual Fisher-density geometry}
\label{sec:rg-spatial-proof-fisher}

\paragraph{Proof of \cref{thm:rg-spatial-fisher-density}.}
Enlarge the chosen natural-parameter compact rectangle to
\(\widehat K=[h_-,h_+]\times[0,j_+]\), with \(h_->0\),
\(0<j_+<\operatorname{atanh}(1/12)\). The positive-\(j\) interior
contains the original compact set. The boundary \(j=0\) is used only for
the direct one-sided continuation below. Write
\(t_+=\tanh j_+\), choose \(\theta>0\) with
\(q=12t_+e^\theta<1\), and put
\(C=16q/(1-q)\), \(\rho=\tanh(h_++2j_+)\), \(d=1-\rho>0\).
These constants work for every volume, depth, support pattern, and fixed
fine exterior. Finite supports and references are parameter independent.
The finite density is an exponential family in \((h,j)\); its field
statistic sums free edge signs and its plaquette statistic includes every
touching plaquette with its exterior completion.

\subsection{Localized source cumulants and actual derivatives}
\label{sec:rg-spatial-proof-higher-response}

Let \(U\) be a finite vector of block observables supported on an anchor
set \(A\), and \(V\) another outside its radius \(R\) neighborhood.
Condition on the fine blocks outside that neighborhood. The weighted
Hamming proof compares any two resulting interior rows. If the original
finite whole-block volume is irregular, its inherited exterior is identical
on both sides and contributes zero disagreement. Only the newly revealed
shell contributes. Thus
\begin{equation}
 \operatorname{TV}(\mathcal L(U,V),\mathcal L(U)\otimes\mathcal L(V))
 \le\min\{1,C|A|B^2(R+1)e^{-\theta(R+1)}\}.
 \label{eq:rg-spatial-proof-group-decoupling}
\end{equation}
Each conditional law of \(U\) is within this bound of their mixture;
integrating with the exterior and marginalizing proves the displayed group
comparison. It passes to infinite constrained laws by finite-cylinder
convergence, also at \(j=0\) by the exact product-fiber law.

For \(m\) bounded observables define
\begin{equation}
 \operatorname{cum}(F_1,\ldots,F_m)
 =\sum_{\pi\in\mathcal P_m}(-1)^{|\pi|-1}(|\pi|-1)!
                              \prod_{D\in\pi}\E\prod_{i\in D}F_i.
 \label{eq:rg-spatial-proof-response-cumulant}
\end{equation}
Its individual absolute bound is \(k_m\prod_i\|F_i\|_\infty\), and a
change of joint law by total variation \(v\) changes it by at most
\(q_m v\prod_i\|F_i\|_\infty\), where
\(k_m=\sum_\pi(|\pi|-1)!\) and \(q_m=2\sum_\pi|\pi|!\).
For \(m=2,3,4\), these pairs are \((2,6),(6,26),(26,150)\).
The second estimate telescopes each product of moments, each moment
changing by at most twice total variation. A cumulant of two independent
nonempty groups vanishes: its generating function factorizes and its
logarithm has no mixed coefficient. Expanding that finite logarithm also
proves \cref{eq:rg-spatial-proof-response-cumulant}.

Let a local \(f\) have anchor support \(A\) of size \(a_0\ge1\), and
let \(r=1,2,3\) labeled edge or plaquette sources have base anchors
\(x_1,\ldots,x_r\). Put \(M=\max_s\operatorname{dist}_\infty(x_s,A)\).
On the complete graph of these labeled anchors and a root representing
\(A\), a minimum-weight spanning tree has an edge of length at least
\(M/r\). Removing a longest edge splits the graph into groups with all
cross distances at least that length; otherwise exchanging a shorter
cross edge would reduce the tree. Plaquette supports enlarge anchor sets
by at most one on either side. For \(M\ge4r\), use
\(R=\lfloor M/r\rfloor-3\) in
\cref{eq:rg-spatial-proof-group-decoupling}; the root group has at most
\(a_0+4r\) anchors. For smaller \(M\), use the individual cumulant
bound. With
\begin{equation}
 L_r(a_0)=e^{4\theta}[k_{r+1}+q_{r+1}C(a_0+4r)],\qquad
 S_r(a_0)=\frac{(8a_0)^rL_r(a_0)(2r+1)!}
                  {(1-e^{-\theta/r})^{2r+2}},
 \label{eq:rg-spatial-proof-response-constants}
\end{equation}
this proves a bound
\(L_r(a_0)B^2(M+1)e^{-\theta M/r}\|f\|_\infty\) for each cumulant.
There are at most \(2B^2\) choices of a source of one derivative type at
each anchor, and at most \(a_0(2M+1)^2\) anchors within distance \(M\).
Counting ordered tuples, including repeats, gives
\begin{equation}
 \begin{split}
 \sum_{s_1,\ldots,s_r}|\operatorname{cum}(f,T_{s_1},\ldots,T_{s_r})|
 &\le S_r(a_0)B^{2r+2}\|f\|_\infty,\\
 \sum_{\max\operatorname{dist}(x_s,A)>L}
       |\operatorname{cum}(f,T_{s_1},\ldots,T_{s_r})|
 &\le S_r(a_0)B^{2r+2}(L+2)^{2r+1}
                     e^{-\theta(L+1)/r}\|f\|_\infty.
 \end{split}
 \label{eq:rg-spatial-proof-response-tails}
\end{equation}
Here \(T_s\) is the appropriate edge or plaquette source for its labeled
derivative direction. To sum the bound, use
\((M+1)^k\le k!{M+k\choose k}\) in the geometric series. For the tail,
write \(M=L+1+n\) and bound
\((M+1)^k\le(L+2)^k(n+1)^k\). Exterior constant sources have zero
cumulant and may safely be overcounted.

Finite differentiation of the normalized exponential family gives these
cumulant sums for each derivative of a parameter-independent observable.
For a varying observable, the actual formula is
\begin{equation}
 \partial^\alpha\E_gf_g
 =\sum_{\beta\le\alpha}{\alpha\choose\beta}
   \sum_{s_1,\ldots,s_{|\alpha-\beta|}}
   \operatorname{cum}_g(\partial^\beta f_g,T_{s_1},\ldots,
                                      T_{s_{|\alpha-\beta|}}),
 \label{eq:rg-spatial-proof-moving-response}
\end{equation}
where an empty source list means ordinary expectation. Differentiating
the finite generating function or its normalized finite sum proves all
multiplicities, including repeated directions. For each fixed source
tuple, finite cylinder convergence gives convergence of its moments,
uniformly on \(\widehat K\). Near tuples form a finite set and far tuples
have the uniform tails in \cref{eq:rg-spatial-proof-response-tails}.
First sending the source cutoff to infinity and then the volume proves
uniform convergence of finite derivatives through order three. The sums
are continuous by the same finite-near, uniform-tail argument. Passing
finite coordinate fundamental-theorem-of-calculus identities to the
limit proves that these are actual continuous derivatives, successively
through order three, with agreeing mixed derivatives. At \(j=0\) the
argument is on right-sided coordinate segments and gives actual one-sided
\(C^3\) continuation. No infinite normalizer or null-event probability
has been differentiated.

\subsection{Full local jets and interaction derivatives}
\label{sec:rg-spatial-proof-local-jets}

For a free path its conditional mean given other fine blocks is the
product \(G_g\) of \(B\) hyperbolic tangents with field
\(h+j a_e\), \(|a_e|\le2\), on at most nine anchor cells. For
\(n\le3\), differentiating the product in labeled directions gives
\begin{equation}
 \|\partial^\beta G_g\|_\infty
 \le g_n B^n\rho^{B-n},\qquad g_n=2^n n!,\quad |\beta|=n.
 \label{eq:rg-spatial-proof-product-jets}
\end{equation}
The first three hyperbolic-tangent derivatives have bounds \(1,2,6\);
each \(j\) differentiation adds a factor at most two. There are at most
\(B^n\) assignments to product factors and the product of their
factorials is at most \(n!\). Negative exponents when \(B<n\) only
enlarge the upper bound. Apply
\cref{eq:rg-spatial-proof-moving-response,eq:rg-spatial-proof-response-tails}
to \(G_g\). With
\begin{equation}
 D_n=g_n+\sum_{b'=0}^{n-1}{n\choose b'}g_{b'}S_{n-b'}(9),\qquad
 M_n(B)=D_nB^{2n+2}\rho^{B-n},\quad M_0(B)=\rho^B,
 \label{eq:rg-spatial-proof-mean-jets}
\end{equation}
every actual order-\(n\) free-path mean derivative is bounded by
\(M_n(B)\). These bounds tend to zero with \(B\).

For patterns agreeing through \(Q_R\), choose
\(L=\lfloor R/2\rfloor-1\) when \(R\ge2\). A near source tuple has
joint support within \(Q_{L+1}\); apply the group comparison on its
remaining buffer and count \([2B^2(2L+1)^2]^n\) ordered tuples. Use
\cref{eq:rg-spatial-proof-response-tails} for each of the two far sums.
Writing \(N_n=2^nq_{n+1}C(1+4n)\), the resulting derivative modulus is
\begin{equation}
\begin{split}
 V_n(B,R)=\min\bigl\{2M_n(B),B^{2n+2}\bigl[
 &N_n(2L+1)^{2n}(R-L)e^{-\theta(R-L)}\\
 &+2S_n(1)(L+2)^{2n+1}e^{-\theta(L+1)/n}\bigr]\bigr\}.
\end{split}
 \label{eq:rg-spatial-proof-jet-modulus}
\end{equation}
For \(R=0,1\) take \(V_n=2M_n\). For values use
\(V_0=\min\{2\rho^B,2CB^2(R+1)e^{-\theta(R+1)}\}\).
The same proof compares a finite whole-block conditional calculation
against its infinite kernel: near sources are interior, and finite
crossing sources and infinite far sources satisfy the two separate tail
bounds. For \(R\ge2\), the bracket in
\cref{eq:rg-spatial-proof-jet-modulus} is at most
\(Z_n(R+1)^{2n+1}e^{-\theta R/(2n)}\), with
\(Z_n=N_n+2e^{\theta/(2n)}S_n(1)\).

Every Bernoulli factor is \((1+s m)/2\), with \(1+s m\ge d\).
Revealing both target parities sequentially retains all their dependence.
Their positive four-state row is actual \(C^3\). For one Bernoulli log
factor its derivative amplitudes are
\begin{equation}
 A_1=\frac{M_1}{d},\qquad
 A_2=\frac{M_2}{d}+\frac{M_1^2}{d^2},\qquad
 A_3=\frac{M_3}{d}+\frac{3M_2M_1}{d^2}+\frac{2M_1^3}{d^3}.
 \label{eq:rg-spatial-proof-log-jet-amplitudes}
\end{equation}
The multiplicities follow by differentiating \(\log(1+s m)\) three
times. Its derivative moduli are
\begin{equation}
 \begin{split}
 T_1={}&V_1/d+M_1V_0/d^2,\\
 T_2={}&V_2/d+(M_2V_0+2M_1V_1)/d^2+2M_1^2V_0/d^3,\\
 T_3={}&V_3/d+(M_3V_0+3M_2V_1+3M_1V_2)/d^2\\
       &+(6M_2M_1V_0+6M_1^2V_1)/d^3+6M_1^3V_0/d^4.
 \end{split}
 \label{eq:rg-spatial-proof-log-jet-moduli}
\end{equation}
Subtract the actual derivative ratios and use
\(|u^{-k}-v^{-k}|\le k|u-v|/d^{k+1}\) to obtain these formulas.
One-cell log probabilities use \(2A_n,2T_n\); reference-state log odds
use \(4A_n,4T_n\). Each modulus can also be capped by twice its
amplitude. The same already revealed first parity is fixed when comparing
the second factor, so the support patterns still agree in their common
interior.

Here is an explicit summability check for the complete interaction jets.
Set
\(E_1=D_1/d\), \(E_2=D_2/d+D_1^2/d^2\), and
\(E_3=D_3/d+3D_2D_1/d^2+2D_1^3/d^3\). Then
\(F_n(B)=4A_n(B)\le4E_nB^{4n}\rho^{B-n}\).
Each \(M_j\) is bounded over all \(B\) by
\(U_j=D_j\rho^{1-j}(2j+2)!/(1-\rho)^{2j+3}\).
Put \(Z_0=2C\) and obtain \(P_n\) from the right side of
\cref{eq:rg-spatial-proof-log-jet-moduli} by replacing every \(M_j\)
with \(U_j\) and every \(V_j\) with \(Z_j\). Thus these are explicit
positive constants; the powers of \(d\) are unchanged. For \(R\ge2\),
\(T_n\le P_nB^{2n+2}(R+1)^{2n+1}e^{-\theta R/(2n)}\).
The constant
\begin{equation}
 H_n=4P_n+
 8E_n\rho^{-n}\frac{(2n-2)!}{(1-\rho)^{2n-1}}e^{\theta/(2n)}
 \label{eq:rg-spatial-proof-jet-shell-constant}
\end{equation}
also covers shell radii one and two using the amplitude cap. Therefore
its order-\(n\) canonical shell is bounded by both \(2F_n(B)\) and
\(H_nB^{2n+2}r^{2n+1}e^{-\theta(r-1)/(2n)}\).
Choose one common \(\lambda_f=\theta/48\). Geometric interpolation,
\(\operatorname{diam}A_{x,r}\le2r\), and the translate count give
\begin{equation}
 \|\partial^\alpha\Phi^{(B)}\|_{\lambda_f}
 \le F_n(B)+\sqrt{2F_n(B)H_n}B^{n+1}e^{\theta/(4n)}
       \sum_{r\ge1}(2r+1)^2r^{n+1}
          e^{-[\theta/(4n)-2\lambda_f]r},\quad |\alpha|=n\le3.
 \label{eq:rg-spatial-proof-complete-jet-norm}
\end{equation}
The bracketed exponent is positive. The positive series is bounded by
\(9(n+3)!/(1-e^{-[\theta/(4n)-2\lambda_f]})^{n+4}\), and its
omitted tail is an explicit convergent bound. The right side decays as a
constant times \(B^{3n+1}\rho^{(B-n)/2}+B^{4n}\rho^{B-n}\).
Values use the preceding canonical-interaction proof in the same weaker
\(\lambda_f\) norm. Finite-shell maps are \(C^3\); their derivatives
through order three converge uniformly in this Banach norm. Successive
coordinate FTC identities prove actual \(C^3\) of the entire interaction.
No shell or derivative has been replaced by a finite truncation in the
statement.

\subsection{Moving expectations and actual marginal information}
\label{sec:rg-spatial-proof-moving-fisher}

For any finite or infinite subset \(S\) of the lexicographic past, let
\(p_S(g,s;y)\) be the four-state predictive row with exactly those past
coordinates prescribed and the other paths, including the future, free.
Let \(l=\log p_S\), with subscripts denoting natural derivatives. Define
\(\mathsf I_{S,ij}=\sum_sp_Sl_il_j\), an observable of the prescribed
past. Its actual first and second parameter derivatives, keeping \(y\)
fixed, are
\begin{equation}
 \begin{split}
 \mathsf I_{S,ij,k}
 &=\sum_sp_S(l_kl_il_j+l_{ik}l_j+l_il_{jk}),\\
 \mathsf I_{S,ij,kv}
 &=\sum_sp_S\bigl[(l_{kv}+l_kl_v)l_il_j
 +l_k(l_{iv}l_j+l_il_{jv})+l_v(l_{ik}l_j+l_il_{jk})\\
 &+l_{ikv}l_j+l_{ik}l_{jv}+l_{iv}l_{jk}+l_il_{jkv}\bigr].
 \end{split}
 \label{eq:rg-spatial-proof-fisher-observable-jets}
\end{equation}
These formulas differentiate the probabilities as well as the scores.
Put \(L_i=2A_i\),
\begin{equation}
 E_B^F=\max\{L_1^2,L_1^3+2L_1L_2,
                L_1^4+5L_1^2L_2+2L_1L_3+2L_2^2\},\qquad
 H_B^F=\max(1,L_1,L_2,L_3).
 \label{eq:rg-spatial-proof-fisher-observable-bound}
\end{equation}
Every matrix entry and its first two derivatives are bounded by
\(E_B^F\). For patterns agreeing in \(Q_R\), every row probability
changes by at most \(V_0\), and each log jet by at most \(2T_n\).
Telescoping the at most ten expanded monomials in
\cref{eq:rg-spatial-proof-fisher-observable-jets}, at most five factors
per term and four states, gives the common modulus
\begin{equation}
 \omega_B(R)=\min\{2E_B^F,
 200(H_B^F)^4[V_0(B,R)+2T_1(B,R)+2T_2(B,R)+2T_3(B,R)]\}.
 \label{eq:rg-spatial-proof-fisher-observable-modulus}
\end{equation}
At fixed \(B\) it is bounded by an explicit constant times
\((R+1)^7e^{-\theta R/6}\). These are differences of actual jets,
not derivatives of an earlier scalar continuity bound.

Expectations now use the original unconstrained fine law, with coarse
observables pulled back by the fixed path map. For a local \(C^2\)
observable on \(Q_N\), put \(a_N=(2N+1)^2\) and
\(W_B(N)=1+2S_1(a_N)B^4+S_2(a_N)B^6\). The moving response formula gives
\begin{equation}
 \max_{|\alpha|\le2}\sup_{g\in\widehat K}
      |\partial^\alpha\E_{\mu_g}f_N(g)|
 \le W_B(N)\max_{|\alpha|\le2}\sup_{g,y}|\partial^\alpha f_N(g,y)|.
 \label{eq:rg-spatial-proof-varying-expectation-bound}
\end{equation}
In particular, the second derivative includes \(\E f_{ij}\), both
source covariance sums involving \(f_i,f_j\), and the complete
three-variable cumulant sum involving \(f\). Formula
\cref{eq:rg-spatial-proof-response-tails} justifies each passage from
finite volumes; neither the observable nor the probability law is held
fixed incorrectly. The function \(W_B(N)\) is polynomial in \(N\).
For a shell series \(f=\sum_Nd_N\) with
\(\sum_NW_B(N)\|d_N\|_{C^2}<\infty\), finite partial sums and their
expectation derivatives are uniformly Cauchy. Repeated FTC gives actual
\(C^2\) of its expectation and termwise differentiation.

Apply this to \(f_N^S=\mathsf I_{S\cap Q_N}\). Its initial shell is
bounded by \(E_B^F\), and its later shell differences through order two
by \(\omega_B(N-1)\). Thus
\begin{equation}
 Z_B^F=W_B(0)E_B^F+\sum_{N\ge1}W_B(N)\omega_B(N-1)<\infty,
 \qquad
 \epsilon_B(R)=2\sum_{N>R}W_B(N)\omega_B(N-1)\longrightarrow0.
 \label{eq:rg-spatial-proof-fisher-expectation-tails}
\end{equation}
If two past sets agree inside \(Q_R\), their finite prefixes agree
exactly through that radius, so their full-law Fisher-observable
expectations differ in entrywise \(C^2\) norm by at most
\(\epsilon_B(R)\). This holds uniformly over the remote shape of the
past. In particular the full-past expectation
\(\mathcal I_B=\E\mathsf I_{\mathrm{past}}\) is actual \(C^2\).

The finite positive marginal on \(\Lambda_L\) factors into its actual
predictive rows in lexicographic order. Each score increment is centered
conditionally on its earlier revealed cells. Two increments are
orthogonal because the earlier one is measurable for the later past.
Consequently the exact finite Fisher matrix is the sum of expected
row-Fisher matrices with their actual finite pasts. Each is a \(C^2\)
function by the preceding varying-expectation argument. Translation
invariance removes the shifts. For \(L\ge2R\), all but at most
\(4RL\) cells have radius \(R\) inside the square; their finite pasts
and full past agree on that radius. Each remaining difference is bounded
by \(2Z_B^F\). Therefore
\begin{equation}
 \|L^{-2}I_{B,\Lambda_L}-\mathcal I_B\|_{C^2,\mathrm{entry}}
 \le\epsilon_B(R)+8Z_B^F R/L.
 \label{eq:rg-spatial-proof-actual-fisher-limit}
\end{equation}
First send \(L\) to infinity at fixed \(R\), then \(R\) to infinity.
This proves the asserted actual marginal \(C^2\) limit and includes all
boundary-score derivatives. The canonical row with future fixed to vacuum
was never substituted for the predictive row.

\subsection{An exhibited positive chart and its curvature}
\label{sec:rg-spatial-proof-regular-chart}

At \(j=0\), with any fixed complete past, the blocks are independent
product parity fibers and the two target path parities are independent
with mean \(b^B\). For one fine edge of a target path set
\begin{equation}
 a(s)=\E[z_e\mid Y=s]-b
      =\frac{s b^{B-1}(1-b^2)}{1+s b^B},\qquad
 v_B=\E a^2=\frac{b^{2B-2}(1-b^2)^2}{1-b^{2B}}>0.
 \label{eq:rg-spatial-proof-product-score-contrast}
\end{equation}
Multiplying the product law by its parity indicator and summing proves
this formula; \(\E a=0\), and conditional-variance decomposition gives
\(v_B\le1-b^2\le1\). The actual field score is
\(s_h=B(a_1+a_2)\). The actual plaquette score is
\(s_j=b^2a_1a_2+\alpha_1a_1+\alpha_2a_2\), with
\(|\alpha_i|\le2B\). To see the mixed term, the plaquette at the common
starting vertex is the unique one meeting both outgoing target paths.
Its other two edges are future paths at \(B=1\), or unused singleton
edges at \(B>1\), and have mean \(b\). Its centered contribution is
\(b^2a_1a_2+b^3(a_1+a_2)\). Every other incident plaquette meets only
one target path, giving a multiple of \(a_i\) of magnitude at most one;
there are \(2B-1\) such incidences. Nonincident plaquettes contribute
zero by product-block independence. These are derivatives of the full
family through its source formula and one-sided limit, not derivatives
of a formula restricted to \(j=0\).

Independence and centering make \(a_1a_2\) orthogonal to both linear
contrasts. The conditional Fisher entries are therefore
\(I_{hh}=2B^2v_B\),
\(I_{hj}=Bv_B(\alpha_1+\alpha_2)\), and
\(I_{jj}=v_B(\alpha_1^2+\alpha_2^2)+b^4v_B^2\).
Its determinant is
\(B^2v_B^2(\alpha_1-\alpha_2)^2+2B^2b^4v_B^3\), while its trace is
at most \(11B^2v_B\). Dividing the determinant by the trace bounds its
smaller eigenvalue below by \(2b^4v_B^2/11\), uniformly in the past.
On a compact positive-field interval with \(b\in[b_-,b_+]\), this is
at least
\begin{equation}
 \eta_{0,B}=\frac2{11}b_-^{4B}(1-b_+^2)^4>0.
 \label{eq:rg-spatial-proof-fisher-product-floor}
\end{equation}

Use the bounds \(A_1,A_2\) on the enlarged rectangle and put
\(Q_B=16(A_1^3+A_1A_2)>0\). The score norm is at most
\(2\sqrt2 A_1\), and its \(j\) derivative at most
\(2\sqrt2 A_2\). Differentiating the four-state Fisher sum gives
operator-norm derivative at most \(Q_B\), including the probability
derivative term. Actual one-sided \(C^3\) and FTC imply
\begin{equation}
 \mathcal I_B(h,j)\succeq\frac{\eta_{0,B}}2 I_2
 \quad\hbox{for}\quad
 0<j<\delta_B:=\min\{j_+/2,\eta_{0,B}/(2Q_B)\}.
 \label{eq:rg-spatial-proof-regular-fisher-region}
\end{equation}
The same floor holds rowwise and hence after the full-law expectation.
An enlarged field interval supplies an open neighborhood of each original
compact interval. For a prescribed finite list of depths, take the finite
minima of positive widths and floors. No positive minimum over every
integer depth is asserted.

On a compact subchart of this region write \(g_{ij}=\mathcal I_{B,ij}\).
The inverse is \(C^2\), by differentiating \(gg^{-1}=I\) twice with its
uniform positive floor. Define
\begin{equation}
 \Gamma^l_{jk}=\frac12 g^{lm}
        (\partial_jg_{mk}+\partial_kg_{mj}-\partial_mg_{jk}),\qquad
 R^l{}_{kij}=\partial_i\Gamma^l_{jk}-\partial_j\Gamma^l_{ik}
             +\Gamma^l_{im}\Gamma^m_{jk}-\Gamma^l_{jm}\Gamma^m_{ik}.
 \label{eq:rg-spatial-proof-fisher-connection-curvature}
\end{equation}
Direct substitution proves torsion freedom and metric compatibility.
Adding and subtracting their cyclic index equations recovers the displayed
\(\Gamma\), proving uniqueness. The covariant-derivative commutator
gives the displayed tensor, so it is coordinate covariant. The connection
is \(C^1\), and curvature, Ricci contraction, and scalar contraction are
continuous because they use metric derivatives through order two.
Equation \cref{eq:rg-spatial-proof-actual-fisher-limit} and the positive
floor make \(L^{-2}I_{B,\Lambda_L}\) invertible for every sufficiently
large \(L\). Continuous inversion and the displayed formulas give compact
convergence of its connections and curvature, at this fixed \(B\).

These are also actual local law charts. At a chart point choose a
sufficiently large finite marginal with positive Fisher matrix. Its
probability Jacobian has rank two because
\(I=(Dp)^\top\operatorname{diag}(1/p)Dp\); hence two atom probabilities
have nonsingular Jacobian. After multiplication by its inverse at the
point, shrink a convex parameter ball until that map's derivative is
within \(1/2\) of the identity. Segment integration proves a lower
Lipschitz bound. The iteration \(x\mapsto x+y-F(x)\) is a contraction
on a smaller ball for nearby \(y\), giving a local inverse.
Differentiating \(F(F^{-1}(y))=y\) gives its derivatives through order
three. This exhibits a chart of the actual law without assuming a global
quotient manifold. \hfill\(\square\)

A vanishing sequence of Fisher tensors does not establish vanishing Fisher
curvature. If a regular metric is multiplied by a positive constant
\(a\), its connection and its Riemann tensor with one raised index are
unchanged, whereas its scalar curvature is divided by \(a\). Thus metric
collapse can coexist with divergent scalar curvature. This elementary
scaling obstruction is separate from the positive fixed-depth charts
proved above. \status{ESTABLISHED}